\packageentry{equations}{amsmath -- Equations and matrices}
\partlabel{Preamble}
\begin{referencecode}
\usepackage{amsmath}
\end{referencecode}
\partlabel{Aligned equations -- one number per line}
\begin{referencecode}
\begin{align}
    F &= ma \label{eq:dynamics} \\
    a &= \frac{F}{m} \label{eq:acceleration}
\end{align}
\end{referencecode}

\begin{align}
    F &= ma \label{eq:dynamics-functional} \\
    a &= \frac{F}{m} \label{eq:acceleration-functional}
\end{align}

\partlabel{Several lines -- one equation number}
\begin{referencecode}
\begin{equation}
    \begin{split}
        E &= T + V \\
          &= \frac{1}{2}mv^2 + mgh
    \end{split}
    \label{eq:energy}
\end{equation}
\end{referencecode}

\begin{equation}
    \begin{split}
        E &= T + V \\
          &= \frac{1}{2}mv^2 + mgh
    \end{split}
    \label{eq:energy-functional}
\end{equation}

Use \texttt{align*} for no numbering, \texttt{\string\notag} to suppress
one line number, and \texttt{\string\eqref\{label\}} for a reference.
\texttt{\&} marks alignment; \texttt{\\} separates lines.

\begin{nobreaks}
\subsubsection*{Matrices and piecewise functions}
\phantomsection\label{matrices}
\begin{referencecode}
\[
    A = \begin{bmatrix}
        1 & 2 \\
        3 & 4
    \end{bmatrix}
    \qquad
    \det(A) = \begin{vmatrix}
        1 & 2 \\
        3 & 4
    \end{vmatrix} = -2
\]
\end{referencecode}

\[
    A = \begin{bmatrix}
        1 & 2 \\
        3 & 4
    \end{bmatrix}
    \qquad
    \det(A) = \begin{vmatrix}
        1 & 2 \\
        3 & 4
    \end{vmatrix} = -2
\]

\texttt{pmatrix}: parentheses; \texttt{bmatrix}: brackets;
\texttt{vmatrix}: determinant bars; \texttt{matrix}: no delimiters.
\end{nobreaks}

\begin{referencecode}
\[
    f(x) = \begin{cases}
        x^2, & \text{if } x \geq 0, \\
        -x,  & \text{if } x < 0.
    \end{cases}
\]
\end{referencecode}

\[
    f(x) = \begin{cases}
        x^2, & \text{if } x \geq 0, \\
        -x,  & \text{if } x < 0.
    \end{cases}
\]

\begin{referencecode}
\[
    \left(\frac{a+b}{c}\right)^2
    \qquad \sin(\theta)
    \qquad \mathbf{v} = (v_x,v_y)
\]
\end{referencecode}

\[
    \left(\frac{a+b}{c}\right)^2
    \qquad \sin(\theta)
    \qquad \mathbf{v} = (v_x,v_y)
\]

Use \texttt{\string\text\{...\}} for words and commands such as
\texttt{\string\sin} or \texttt{\string\log} for operators.

